\documentclass[11pt,letterpaper]{article}
\usepackage[T1]{fontenc}
\usepackage[utf8]{inputenc}
\usepackage{lmodern}
\usepackage{microtype}
\usepackage[margin=1in,headheight=14pt]{geometry}
\usepackage{amsmath,amssymb,amsthm,mathtools}
\usepackage{enumitem,booktabs,array,longtable}
\usepackage{needspace,etoolbox}
\usepackage[dvipsnames]{xcolor}
\usepackage{hyperref}
\usepackage{fancyhdr}
\usepackage{listings}
\hypersetup{colorlinks=true,linkcolor=MidnightBlue,citecolor=MidnightBlue,urlcolor=MidnightBlue,
 pdftitle={Small Divisors at Surreal Scales},pdfauthor={Research draft prepared with ChatGPT}}
\pagestyle{fancy}\fancyhf{}
\fancyhead[L]{\small Small divisors at surreal scales}
\fancyhead[R]{\small Research draft}
\fancyfoot[C]{\thepage}
\setlength{\parskip}{0.35em}
\setlength{\parindent}{1.3em}
\setlist{topsep=4pt,itemsep=2pt,parsep=0pt}
\newtheorem{theorem}{Theorem}[section]
\newtheorem{lemma}[theorem]{Lemma}
\newtheorem{proposition}[theorem]{Proposition}
\newtheorem{corollary}[theorem]{Corollary}
\theoremstyle{definition}
\newtheorem{definition}[theorem]{Definition}
\newtheorem{example}[theorem]{Example}
\theoremstyle{remark}
\newtheorem{remark}[theorem]{Remark}
\newcommand{\No}{\mathbf{No}}
\newcommand{\R}{\mathbb R}
\newcommand{\C}{\mathbb C}
\newcommand{\Z}{\mathbb Z}
\newcommand{\N}{\mathbb N}
\newcommand{\Q}{\mathbb Q}
\newcommand{\T}{\mathbb T}
\newcommand{\A}{\mathcal A}
\newcommand{\HH}{\mathcal H}
\newcommand{\supp}{\operatorname{supp}}
\newcommand{\rank}{\operatorname{rank}}
\newcommand{\id}{\operatorname{id}}
\newcommand{\im}{\operatorname{Im}}
\newcommand{\ord}{\operatorname{ord}}
\newcommand{\diver}{\operatorname{div}}
\newcommand{\mean}[1]{\langle #1\rangle}
\newcommand{\mon}[1]{\langle #1\rangle_{\N}}
\newcommand{\Lom}{L_{\Omega}}
\newcommand{\Gom}{G_{\Omega}}
\newcommand{\res}{\mathcal N_{\Omega}}
\newcommand{\SDcond}{\mathrm{SS}}
\newcommand{\coef}[2]{[t^{#1}]\,#2}
\newcommand{\valuation}{v}
\newcommand{\eps}{\varepsilon}
\newcommand{\dd}{\,\mathrm d}
\DeclareMathOperator*{\limsupk}{lim\,sup}
\numberwithin{equation}{section}
\allowdisplaybreaks[2]
\BeforeBeginEnvironment{theorem}{\Needspace{9\baselineskip}}
\BeforeBeginEnvironment{lemma}{\Needspace{5\baselineskip}}
\BeforeBeginEnvironment{proposition}{\Needspace{6\baselineskip}}
\BeforeBeginEnvironment{corollary}{\Needspace{5\baselineskip}}
\begin{document}

\begin{titlepage}
\centering
\vspace*{1.0cm}
{\large SURREAL AND SURCOMPLEX ANALYSIS\par}
\vspace{1.1cm}
{\LARGE\bfseries Small Divisors at Surreal Scales\par}
\vspace{0.6cm}
{\Large Resonance flags, analytic lifting obstructions,\par
and Hahn normal forms\par}
\vspace{1cm}
{\large A research article with proofs and exact finite checks\par}
\vspace{0.5cm}
{September 21, 2026\par}
\vspace{1cm}
\begin{minipage}{0.89\textwidth}
\begin{abstract}
A surreal frequency can contain a transfinite hierarchy of scales, while a
cohomological equation on a torus has infinitely many Fourier divisors.
We show that a frequency vector in $\No^d$ nevertheless has at most $d$
nonzero divisor-valuation levels.  A canonical descending flag of integer
resonance lattices identifies these levels and gives a single well-ordered
support containing every reciprocal divisor.

For holomorphic coefficients on one fixed complex strip, a stratified
subexponential small-divisor condition is then necessary and sufficient for
universal Hahn-analytic solvability.  Under the same condition, a vector field
perturbed beyond its last resonance scale has a unique infinitesimal,
mean-zero conjugacy to a constant vector field with a frequency correction.
The proof uses finite word representations in positive Hahn supports, not an
invalid appeal to sequential convergence.  Conversely, failure of the
arithmetic condition produces non-linearizable perturbations beyond any
prescribed valuation threshold.

We also construct entire forcing terms for which any prescribed finite
number of analytic lifting steps succeeds, but the next forced coefficient
is not even a distribution.  A boundary example makes the perturbation-scale
threshold sharp as a uniform statement.  An invariant-density theorem and
explicit examples at scales $1,t,t^\omega$ complete the development.
\end{abstract}
\end{minipage}
\vfill
\begin{minipage}{0.89\textwidth}
\small
\textbf{Status.} These are proposed research results with detailed proofs,
not a claim to have settled a named longstanding conjecture.  The precise
combination of results was not located in the focused literature review;
priority is not guaranteed.  Classical inputs are identified explicitly.
The article is unrefereed and has not been formalized in a proof assistant.
The accompanying program verifies finite identities only.
\end{minipage}
\end{titlepage}

{\small\setlength{\parskip}{0pt}\tableofcontents}
\newpage

\section{The question, the answer, and the scope}
\label{sec:intro}

Consider the equation
\begin{equation}\label{eq:cohom-intro}
 \Lom u=f,\qquad
 \Lom=\Omega_1\frac{\partial}{\partial\theta_1}
       +\cdots+\Omega_d\frac{\partial}{\partial\theta_d},
 \qquad \theta\in(\R/2\pi\Z)^d.
\end{equation}
For a real frequency vector, Fourier transformation asks us to divide by
$i\,k\cdot\Omega$ for each integer vector $k$.  With surreal frequencies,
there are two independent obstacles.  Divisors can be small as ordinary real
numbers, destroying Fourier decay; and their reciprocals can have complicated
Hahn supports, potentially destroying the existence of a single surreal-valued
answer.  Solving the scalar equation separately at each Fourier mode does not,
by itself, solve either problem.

The main outcome is a finite reduction of this apparently transfinite
problem.  There are finitely many leading real forms, each acting on a
successively smaller integer lattice.  They control all small-divisor
arithmetic.  The remaining surreal coefficients affect the support and the
higher coefficient multipliers, but introduce no additional arithmetic test.
That reduction also supports a non-linear conjugacy theorem.

\subsection{Main statements in plain mathematical terms}

Every result is proved in a set-sized Hahn workspace; the connection to the
full surreal field is explained in Section~\ref{sec:workspace}.
The four central conclusions are as follows.

\begin{enumerate}[label=\textbf{\arabic*.},leftmargin=*]
\item A vector with $d$ surreal components has at most $d$ distinct values of
$v(k\cdot\Omega)$ among its nonzero integer divisors.  Their reciprocal
supports lie in one explicitly constructed well-ordered set
(Theorem~\ref{thm:flag}).
\item The equation~\eqref{eq:cohom-intro} is universally solvable in a fixed-strip
Hahn-analytic algebra exactly when the leading real divisors on every
resonance stratum have subexponential reciprocals
(Theorem~\ref{thm:linear}).
\item The same condition gives an exact infinitesimal conjugacy for every
perturbation beyond the last resonance scale.  It is also necessary for
universal such linearization, even when perturbations may be required to have
arbitrarily high valuation (Theorems~\ref{thm:nonlinear}
and~\ref{thm:nonlinear-necessity}).
\item Analytic finite jets alone cannot replace the arithmetic condition.
There is an explicit family in which the first $N+1$ solution coefficients are
analytic but the next coefficient is not a distribution, for arbitrary $N$
(Theorem~\ref{thm:jet}).
\end{enumerate}

The non-linear theorem is more than a calculation of a few perturbative
coefficients: it constructs a Hahn sum with possibly transfinite support,
proves a common analyticity domain for all its coefficients, proves uniqueness,
and identifies the exact scale beyond which a uniform theorem holds.
It is nevertheless a coefficientwise Hahn theorem, \emph{not} a convergence
theorem for a real perturbation parameter.

\subsection{Relationship to the supplied repository}

The repository catalogue and root documentation were inspected at commit
\texttt{aa846271b4dc}~\cite{repo}.
They already describe substantial work on Hahn evaluation, common-domain
analytic geometry, finite deformations, global divisors, polynomial algebra,
trigonometry, differential equations, rank-one analytic geometry, and finite
spectral theory.  The present article does not repackage those subjects.
Its focus is the interaction of \emph{infinitely many Fourier modes} with
\emph{several surreal divisor scales}.  A targeted repository code search for
``cohomological'' returned no matches; this is not an exhaustive audit of every
manuscript or every synonym.

Two distinctions in that catalogue are especially relevant here.  First,
Hahn summation is not convergence in the fine topology of the whole surreal
field.  Second, a separate ordinary analytic germ for each Hahn coefficient
need not give a common domain.  We use one fixed open strip throughout and
supply global support certificates.  Both requirements are part of each main
theorem, not implicit conventions.

\subsection{What is classical, and what is proposed here}

The scalar Fourier division mechanism and its ordinary subexponential
arithmetic criterion belong to classical small-divisor theory.  Bounemoura
provides a systematic discussion of linear and non-linear arithmetic
conditions~\cite{bounemoura}; Kamie\'nski studies detailed analytic solvability
and Liouville obstructions for a two-variable cohomological
operator~\cite{kamen}.  We do not claim either ingredient as new.
Hahn inverses and positive-support summability are also classical
\cite{neumann,poonen}.

The proposed contributions are the finite resonance-flag formulation for
arbitrary surreal frequencies, the mode-uniform Hahn support theorem, their
exact fixed-strip solvability consequence, the explicit arbitrary-order
lifting obstruction, and the support-controlled non-linear theorem with its
necessity and sharp-scale statements.  Each proof is given below, rather than
being inferred from a literature search.

For comparison, the 2026 preprint of Chevallier, Lopes Dias, Gaiv\~ao, Marnat,
and Moshchevitin treats quantitative \emph{ordinary analytic} linearization
through Brjuno arithmetic~\cite{chevallier}.  Our theorem neither strengthens
nor contradicts that kind of convergence result.  No bound is imposed here on
the growth of analytic seminorms as the Hahn exponent increases.  Section
\ref{sec:limits} explains exactly why this difference matters.

\section{Set-sized workspaces and fixed-strip Hahn functions}
\label{sec:workspace}

\subsection{The scalar field and the sign convention}

Let $\Gamma$ be a set-sized ordered abelian group.  We write
\[
 K_\Gamma=\C((t^\Gamma)),\qquad
 F_\Gamma=\R((t^\Gamma)).
\]
An element is a formal sum $a=\sum_{\gamma\in S}a_\gamma t^\gamma$ with
$S\subseteq\Gamma$ well ordered in the increasing order.  For $a\ne0$, put
$v(a)=\min\supp(a)$, and put $v(0)=+\infty$.  Thus larger valuation means a
smaller scale.

To obtain literal surreal and surcomplex numbers, take
$\Gamma\subseteq(\No,+)$ and interpret
\begin{equation}\label{eq:omega-map}
 t^\gamma=\omega^{-\gamma}.
\end{equation}
The right-hand side uses Conway's normal-form omega-map.  It is not an
assertion that this map is Gonshor's exponential.  Normal forms give the
usual embedding of $F_\Gamma$ in $\No$, and adjoining $i$ embeds $K_\Gamma$ in
$\No[i]$; see Gonshor~\cite{gonshor}.  In this interpretation, positive-valuation
scalars are infinitesimal.

Any finite vector of surreal numbers is covered: collect all exponents in
its component normal forms and take the additive subgroup they generate.
This is a set because there are only finitely many set-sized supports.
Additional set-sized families of data can be included by taking another
set-sized union.  No set of all surreal numbers and no proper-class Hahn
support is needed.

The proofs below use only ordered-group Hahn algebra.  Consequently, they
remain valid for abstract $\Gamma$ and commute with order-preserving
extensions of the value group.  Divisibility of $\Gamma$ is not required.

\subsection{Analytic coefficients on a common strip}

Fix $\rho>0$ and set
\[
 \T^d_\rho=\{z\in\C^d/(2\pi\Z)^d:
                 |\im z_j|<\rho\text{ for all }j\},
 \qquad \A_\rho=\mathcal O(\T^d_\rho).
\]
Functions are periodic and holomorphic on this \emph{open} strip.  They need
not be bounded at its boundary.  With
$|k|=|k_1|+\cdots+|k_d|$, write
\[
 f(\theta)=\sum_{k\in\Z^d}\widehat f(k)e^{ik\cdot\theta},
 \qquad
 \|f\|_r=\sum_{k\in\Z^d}|\widehat f(k)|e^{r|k|}
 \quad(0<r<\rho).
\]
A periodic holomorphic function belongs to $\A_\rho$ precisely when all these
seminorms are finite.  Indeed, boundedness on any narrower closed strip and
Cauchy's Fourier estimate give exponential decay there; inserting an
intermediate strip gives the weighted $\ell^1$ estimate.  Conversely, these
estimates give locally uniform convergence of the Fourier series and its
termwise derivatives on every narrower strip.

In particular, $\A_\rho$ is closed under finite sums, products, and all finite
orders of differentiation.  Fourier projection onto \emph{any} subset of
$\Z^d$ preserves $\A_\rho$, since it only removes terms from $\|f\|_r$.

\begin{lemma}[Subexponential multipliers]\label{lem:multiplier}
Suppose $m:\Z^d\to\C$ satisfies, for every $\eps>0$,
\[
 |m(k)|\le C_\eps e^{\eps|k|}.
\]
Then the Fourier multiplier $\widehat{Tf}(k)=m(k)\widehat f(k)$ maps
$\A_\rho$ to $\A_\rho$.  Every finite product of such multipliers and
polynomial functions of $k$ has the same property.
\end{lemma}
\begin{proof}
Given $r<r'<\rho$, use the estimate for $\eps=r'-r$ to obtain
$\|Tf\|_r\le C_{r'-r}\|f\|_{r'}$.  A fixed polynomial is bounded by a constant
times $e^{\eta|k|}$ for every $\eta>0$.  Splitting an arbitrary exponent
budget among finitely many factors proves the last assertion.
\end{proof}

\begin{definition}[The common-strip Hahn algebra]
\[
 \HH_\rho=\A_\rho((t^\Gamma))
 =\left\{f=\sum_{\beta\in B}t^\beta f_\beta:
          B\text{ well ordered},\ f_\beta\in\A_\rho\right\}.
\]
The support is the set of exponents whose coefficient function is nonzero.
For a vector, the support is the union of the component supports and its
valuation is their minimum.  The zero vector has valuation $+\infty$.
\end{definition}

Write $\HH_{\rho,>0}=\{f\in\HH_\rho:v(f)>0\}$; this notation includes zero.
No bound on $\|f_\beta\|_r$ uniform in $\beta$ is imposed.  This is an
important deliberate choice.  On the other hand, the same $\rho$ serves
every $\beta$; independent coefficient germs do not suffice.

Fourier coefficients and torus mean are defined coefficientwise:
\[
 \widehat f(k)=\sum_\beta t^\beta\widehat{f_\beta}(k)\in K_\Gamma,
 \qquad
 \mean f=\widehat f(0).
\]
The coefficient field acts on $\HH_\rho$, and ordinary torus derivatives act
on each $f_\beta$.  Fourier coefficients determine an element of $\HH_\rho$
uniquely, by the ordinary uniqueness theorem at each Hahn exponent.

\subsection{Two sums, and two derivatives, that must not be confused}

The Fourier series of an ordinary coefficient is an ordinary complex
analytic sum.  The series in $t^\Gamma$ is a Hahn sum.  A family with
infinitely many nonzero terms of valuation zero need not be Hahn-summable,
even when its complex coefficients form an absolutely convergent series.
Our arguments group all Fourier modes \emph{inside each analytic coefficient}
first.  We never assert that the unrestricted two-index family of Fourier
and Hahn terms is Hahn-summable.

Likewise, $\partial/\partial\theta_j$ is an \emph{external} derivation that
annihilates $K_\Gamma$.  It is not the Berarducci--Mantova derivation on $\No$.
For example, $\partial_\theta\sin\theta=\cos\theta$ here says nothing about
whether a surreal $y$ can solve $\partial^2y+y=0$ for that internal derivation.
This convention keeps the present theory separate from the repository's
internal differential-equation report~\cite{repo}.

\subsection{Actual values on infinitesimal thickenings}

The algebra is not merely a collection of scalar formal symbols.  If
$z_0\in\T^d_\rho$ is ordinary and $\zeta\in K_\Gamma^d$ has positive
valuation, define
\begin{equation}\label{eq:monad-eval}
 f(z_0+\zeta)=
 \sum_{\beta}\sum_{\alpha\in\N^d}
 t^\beta\frac{\partial^\alpha f_\beta(z_0)}{\alpha!}\zeta^\alpha.
\end{equation}
The support lemma in the next section proves that this is a Hahn sum and
hence a surcomplex number under~\eqref{eq:omega-map}.
Thus $\HH_\rho$ describes functions on infinitesimal thickenings of an
ordinary complex torus strip.  We do not claim continuity from an ordinary
torus into the fine topology of the entire $\No[i]$.

\section{The support calculus}
\label{sec:support}

The following classical fact is the algebraic foundation of the paper.
Its positive-support form is often called Neumann's lemma; see
\cite{neumann,poonen}.

\begin{lemma}[Positive-support word lemma]\label{lem:neumann}
If $E\subseteq\Gamma_{>0}$ is well ordered, the additive monoid
\[
 M=\mon E=\{e_1+\cdots+e_n:n\ge0,\ e_j\in E\}
\]
is well ordered.  Every $\gamma\in\Gamma$ has only finitely many
representations as a finite ordered word sum with letters in $E$, including
representations of different lengths.  The empty word represents zero.
Moreover, the sum of finitely many well-ordered subsets of $\Gamma$ is well
ordered, with finitely many decompositions of any fixed element.
\end{lemma}

We use the full assertion, including finiteness across word lengths.  It
implies that if $h$ has positive support, the formal inverse
$(1-h)^{-1}=\sum_{n\ge0}h^n$ is a Hahn identity.  It also justifies
\eqref{eq:monad-eval}, finite-dimensional Taylor substitution, and the marked
word construction in the non-linear proof.

\begin{lemma}[Taylor substitution]\label{lem:taylor}
Let $h\in\HH_\rho^d$ have positive valuation.  For every $f\in\HH_\rho$,
\begin{equation}\label{eq:substitution}
 S_hf=f\circ(\id+h):=
 \sum_\beta t^\beta
 \sum_{\alpha\in\N^d}
 \frac{\partial^\alpha f_\beta}{\alpha!}h^\alpha
\end{equation}
is well defined in $\HH_\rho$.  If $\supp f\subseteq B$ and
$\supp h\subseteq C\subseteq\Gamma_{>0}$, then
\[
 \supp(S_hf)\subseteq B+\mon C.
\]
Taylor substitution is a ring homomorphism, obeys the chain rule, and is
associative whenever the substitutions are infinitesimal.  It satisfies
\begin{align}
 v(S_hf)&\ge v(f),\label{eq:sub-order}\\
 v(S_hf-S_gf)&\ge v(f)+v(h-g)\label{eq:sub-lipschitz}
\end{align}
for positive-valuation $h,g$.
\end{lemma}
\begin{proof}
At each exponent the positive-support word lemma permits only finitely
many Taylor monomials and finitely many decompositions involving $B$.
Every resulting coefficient is a finite sum of products of derivatives of
functions in $\A_\rho$, hence lies in that same algebra.  The formal
Taylor product, composition, and chain-rule identities can therefore be
checked coefficientwise using finite polynomial identities; no limiting
interchange is needed.  For~\eqref{eq:sub-lipschitz}, subtract corresponding
Taylor monomials and telescope their products.  Every term has one factor
from $h-g$, and all its other $h$ or $g$ factors have positive valuation.
The term with multi-index zero cancels.  This gives the claimed bound.
\end{proof}

A particularly dangerous shortcut is worth excluding now.  If
$\Gamma=\Z+\Z\omega$ and $v(h)=1$, then $v(h^n)=n$ never exceeds $\omega$.
Thus the partial sums of $\sum h^n$ need not converge in the intrinsic
valuation topology, let alone in the fine topology of $\No[i]$.  The Hahn
identity is nevertheless valid because each exponent has finitely many word
representations.  The proof in Section~\ref{sec:nonlinear} uses precisely
this stronger support mechanism.

\section{A finite flag for infinitely many surreal divisors}
\label{sec:flag}

Let $\Omega\in F_\Gamma^d\setminus\{0\}$ and collect its component normal
forms into
\begin{equation}\label{eq:frequency}
 \Omega=\sum_{\beta\in S}t^\beta\omega_\beta,
 \qquad \omega_\beta\in\R^d\setminus\{0\},
 \quad S\text{ well ordered}.
\end{equation}
The exact resonance lattice is
\[
 \res=\{k\in\Z^d:k\cdot\Omega=0\}
       =\bigcap_{\beta\in S}\ker(k\mapsto k\cdot\omega_\beta).
\]

\subsection{Construction of the flag}

Start with $\Lambda_0=\Z^d$.  If
$\Lambda_{j-1}\ne\res$, let $\gamma_j$ be the least exponent of $S$ whose
coefficient form does not vanish identically on $\Lambda_{j-1}$.  Put
\begin{equation}\label{eq:flag-def}
 a_j(k)=k\cdot\omega_{\gamma_j},\qquad
 \Lambda_j=\{k\in\Lambda_{j-1}:a_j(k)=0\},\qquad
 \Sigma_j=\Lambda_{j-1}\setminus\Lambda_j.
\end{equation}
The $\Sigma_j$ are the nonresonant strata.  The construction stops as soon
as the current lattice is $\res$.

\begin{theorem}[Finite flag and common reciprocal support]\label{thm:flag}
The construction terminates after $r\le d$ steps and has the following
properties.
\begin{enumerate}[label=(\roman*),leftmargin=*]
\item Each $\Lambda_j$ is a saturated subgroup of $\Z^d$, its rank strictly
decreases at each step, and
\[
 \gamma_1<\cdots<\gamma_r,\qquad \Lambda_r=\res.
\]
The strata partition $\Z^d\setminus\res$.
\item For $k\in\Sigma_j$,
\[
 k\cdot\Omega=t^{\gamma_j}a_j(k)(1+q_k),
 \qquad v(q_k)>0,
 \qquad v(k\cdot\Omega)=\gamma_j,
\]
where $q_k=0$ is permitted.
\item Define
\begin{align}
 E_j&=\{\beta-\gamma_j:\beta\in S,\ \beta>\gamma_j\},\label{eq:E}\\
 E&=\bigcup_{j=1}^rE_j,\qquad M_E=\mon E,\qquad
 T_\Omega=\bigcup_{j=1}^r(-\gamma_j+M_E).\label{eq:T}
\end{align}
Then $E$ is well ordered and positive, $T_\Omega$ is well ordered, and
\[
 \supp\bigl((k\cdot\Omega)^{-1}\bigr)\subseteq T_\Omega
 \quad\text{for every }k\notin\res.
\]
If $\kappa=\gamma_r$, then $T_\Omega\subseteq[-\kappa,+\infty)$.
\end{enumerate}
\end{theorem}
\begin{proof}
Each lattice is the common kernel of real additive homomorphisms on
$\Z^d$, hence is saturated: if a nonzero integer multiple of $k$ lies in the
kernel, so does $k$.  At a nonterminal step $a_j$ is a nonzero homomorphism
on the finitely generated free abelian group $\Lambda_{j-1}$.  Its image is a
nonzero torsion-free group, so the rank of the kernel is strictly smaller.
At most $d$ such decreases are possible.

All coefficient forms below $\gamma_j$ vanish on $\Lambda_{j-1}$ by the
minimal choice of $\gamma_j$.  The form at $\gamma_j$ also vanishes on
$\Lambda_j$.  Therefore the next selected exponent is strictly larger.
Termination can occur only when every coefficient form vanishes on the
current lattice, which is exactly $\res$.  This proves (i), and the same
observation identifies the leading nonzero term for each $k\in\Sigma_j$.
Explicitly,
\[
 q_k=\sum_{\substack{\beta\in S\\\beta>\gamma_j}}
       \frac{k\cdot\omega_\beta}{a_j(k)}t^{\beta-\gamma_j}.
\]
Its support is contained in $E_j$ and is positive, proving (ii).

A tail of a well-ordered set, translated by a fixed element, is well
ordered.  Finite unions of well-ordered subsets of a linear order are well
ordered.  These facts prove the assertion about $E$; Lemma~\ref{lem:neumann}
then proves the assertions about $M_E$ and $T_\Omega$.  Finally,
\begin{equation}\label{eq:reciprocal}
 \frac1{k\cdot\Omega}=
 \frac{t^{-\gamma_j}}{a_j(k)}
 \sum_{n\ge0}(-q_k)^n\qquad(k\in\Sigma_j)
\end{equation}
is a Hahn identity with the asserted support.  Since every element of $M_E$
is nonnegative and $\gamma_j\le\kappa$, the final lower bound follows.
\end{proof}

The number $r$ is a count of \emph{arithmetically visible} scales, not the
number of monomials in $\Omega$.  Coefficients whose forms vanish on a
remaining resonance lattice may be skipped, even across a long transfinite
interval of exponents.  The result does not assert that only $d$ monomials of
$\Omega$ matter to the eventual solution: all later coefficients can affect
higher Hahn coefficients of the inverse.

\begin{proposition}[Coefficient structure of the inverse]\label{prop:coeff-inverse}
Fix $\eta\in\Gamma$ and a stratum $\Sigma_j$.  On that stratum, the function
\[
 k\longmapsto \coef{\eta}{(k\cdot\Omega)^{-1}}
\]
is a finite sum of expressions
\begin{equation}\label{eq:mult-form}
 c\,\frac{\prod_{\ell=1}^n(k\cdot\omega_{\beta_\ell})}
                {a_j(k)^{n+1}},
 \qquad
 \sum_{\ell=1}^n(\beta_\ell-\gamma_j)=\eta+\gamma_j,
 \quad\beta_\ell>\gamma_j,
\end{equation}
where $c$ is a real constant.  Empty products are included when $n=0$.
The number of terms and the denominator powers can depend on $\eta$, but
are finite for each fixed $\eta$.
\end{proposition}
\begin{proof}
Expand~\eqref{eq:reciprocal} and extract the coefficient.  All words have
letters in $E_j$.  Lemma~\ref{lem:neumann} gives finitely many such words of
all lengths.  The map $\beta\mapsto\beta-\gamma_j$ is injective, so this
also bounds the number of coefficient choices.  Each word gives the indicated
product and denominator.
\end{proof}

\begin{remark}[A support bound is not summability over the modes]
The family $((k\cdot\Omega)^{-1})_{k\notin\res}$ is generally not a
Hahn-summable family: infinitely many leading terms may occupy the same
exponent.  Theorem~\ref{thm:flag} supplies a \emph{common support}, not a Hahn
sum over $k$.  Proposition~\ref{prop:coeff-inverse} will allow ordinary Fourier
summation at each coefficient; that is the additional analytic step.
\end{remark}

\begin{proposition}[An arithmetic-free valuation bound]\label{prop:val-bound}
If $h\in\HH_\rho$ has no Fourier modes in $\res$ and $\Lom h=g$, then
\[
 v(h)\ge v(g)-\kappa.
\]
This assertion does not assume any small-divisor estimate.
\end{proposition}
\begin{proof}
For a nonzero $h$, let $\beta=v(h)$.  The nonzero analytic function
$h_\beta$ has a nonzero Fourier coefficient at some $k\notin\res$.
Consequently $v(\widehat h(k))=\beta$, and
\[
 v(\widehat g(k))=\beta+v(k\cdot\Omega)\le\beta+\kappa.
\]
Every Fourier coefficient of $g$ has valuation at least $v(g)$.
Rearranging gives the result.  The zero case is immediate.
\end{proof}

\section{An exact arithmetic criterion for Hahn-analytic solvability}
\label{sec:linear}

\begin{definition}[Stratified subexponential condition]\label{def:SS}
The frequency $\Omega$ satisfies $(\SDcond)$ if, for each flag stratum
$\Sigma_j$ and every real $\eps>0$, there is a real $C_{j,\eps}>0$ such that
\begin{equation}\label{eq:SS}
 \frac1{|a_j(k)|}\le C_{j,\eps}e^{\eps|k|}
 \qquad(k\in\Sigma_j).
\end{equation}
Equivalently, on each stratum,
\[
 \limsup_{\substack{|k|\to\infty\\k\in\Sigma_j}}
 \frac{\log^+(1/|a_j(k)|)}{|k|}=0.
\]
\end{definition}

This condition concerns real leading coefficients, not the surreal size of
the whole divisor.  A divisor can be infinitesimal because its scale is
$t^{\gamma_j}$ and still have completely harmless arithmetic on its stratum.
Conversely, a divisor at scale one can have disastrous real Fourier
arithmetic.  No comparison between a positive real number and a positive
infinitesimal substitutes for~\eqref{eq:SS}.

Let $\Pi_{\res}$ denote coefficientwise Fourier projection onto the exact
resonance lattice.  It is defined on $\HH_\rho$ because arbitrary Fourier
projection preserves $\A_\rho$.

\begin{theorem}[Universal fixed-strip solvability]\label{thm:linear}
For a nonzero frequency $\Omega\in F_\Gamma^d$, the following are equivalent.
\begin{enumerate}[label=(\roman*),leftmargin=*]
\item $\Omega$ satisfies $(\SDcond)$.
\item For every $\rho>0$ and every $f\in\HH_\rho$ with
$\Pi_{\res}f=0$, there exists a unique $u\in\HH_\rho$ satisfying
\[
 \Lom u=f,\qquad \Pi_{\res}u=0.
\]
\end{enumerate}
When these conditions hold, the solution operator $\Gom$ is given by
\begin{equation}\label{eq:green}
 \widehat{\Gom f}(k)=
 \begin{cases}
 \widehat f(k)/(i\,k\cdot\Omega),&k\notin\res,\\
 0,&k\in\res.
 \end{cases}
\end{equation}
It has the support and valuation estimates
\begin{equation}\label{eq:green-support}
 \supp(\Gom f)\subseteq\supp(f)+T_\Omega,
 \qquad v(\Gom f)\ge v(f)-\kappa.
\end{equation}
Without normalization, the kernel consists exactly of Hahn-analytic
functions with Fourier support contained in $\res$.

If $(\SDcond)$ fails, there is already an ordinary real analytic mean-compatible
forcing, supported on one stratum, for which a forced coefficient of any
putative solution is not a distribution on the real torus.
\end{theorem}

\begin{proof}[Proof of sufficiency]
For $\eta\in T_\Omega$, define a scalar Fourier multiplier
\[
 m_\eta(k)=
 \begin{cases}
 \coef{\eta}{(i\,k\cdot\Omega)^{-1}},&k\notin\res,\\
 0,&k\in\res.
 \end{cases}
\]
On each stratum, Proposition~\ref{prop:coeff-inverse} expresses $m_\eta$ as
a finite sum of rational expressions with a polynomial numerator and a
finite power of $a_j(k)$ in the denominator.  Condition $(\SDcond)$ and
Lemma~\ref{lem:multiplier} show that each is a subexponential multiplier.
There are only finitely many strata, and projection onto a stratum does not
affect the estimate.  Hence $m_\eta$ defines an operator
$G_\eta:\A_\rho\to\A_\rho$.

Write $f=\sum_{\beta\in B}t^\beta f_\beta$.  Set
\begin{equation}\label{eq:green-coefficient}
 u_\gamma=\sum_{\substack{\beta\in B,\ \eta\in T_\Omega\\
                         \beta+\eta=\gamma}}G_\eta f_\beta.
\end{equation}
The sum is finite at each $\gamma$, and $B+T_\Omega$ is well ordered.
Therefore $u=\sum_\gamma t^\gamma u_\gamma$ belongs to $\HH_\rho$.
This proves the common-domain and common-support assertions rather than
assuming them from the modewise formula.

Fourier extraction in~\eqref{eq:green-coefficient} now gives
\eqref{eq:green}.  Multiplying at each Fourier mode by $i\,k\cdot\Omega$
proves the equation.  Uniqueness and the kernel description follow by the
same calculation, since $K_\Gamma$ is a field.  The support estimate follows
from the construction; $T_\Omega\ge-\kappa$ gives the valuation estimate.
\end{proof}

\begin{proof}[Proof of necessity]
If $(\SDcond)$ fails on $\Sigma_j$, choose $\eps>0$ and distinct modes $k_n$ on
that stratum, no two equal up to sign, with $|k_n|\to\infty$ and
\[
 |a_j(k_n)|^{-1}>e^{2\eps|k_n|}.
\]
This follows from the failure of the subexponential bound after decreasing
the exponent if necessary.  The real function
\begin{equation}\label{eq:bad-forcing}
 f(\theta)=2\sum_{n\ge1}e^{-\eps|k_n|}\cos(k_n\cdot\theta)
\end{equation}
belongs to $\A_\eps$, has no resonant modes, and is an element of
$\HH_\eps$ supported at Hahn exponent zero.

For any solution, the scalar Fourier equation forces
\[
 \widehat u(k_n)=\frac{e^{-\eps|k_n|}}{i\,k_n\cdot\Omega}.
\]
The coefficient of $t^{-\gamma_j}$ consequently has Fourier coefficients
of modulus greater than $e^{\eps|k_n|}$.  Fourier coefficients of a
distribution on a compact torus have at most polynomial growth: this follows
by applying its finite-order continuity estimate to the Fourier characters.
Thus that forced coefficient is not a distribution, and in particular is
not holomorphic on any positive-width strip.  Adding resonant modes cannot
alter these nonresonant coefficients.  This disproves (ii).
\end{proof}

\subsection{Consequences and quantitative information}

\begin{corollary}[Divisor loss is sharp]\label{cor:sharp-linear}
The valuation loss $\kappa$ in~\eqref{eq:green-support} cannot be improved
uniformly.  For any $k\in\Sigma_r$ and any $s\in\Gamma$, the solution for
$f=t^s e^{ik\cdot\theta}$ has valuation $s-\kappa$.
\end{corollary}
\begin{proof}
Its unique nonzero Fourier mode is divided by a scalar of valuation
$\kappa$.
\end{proof}

A convenient sufficient arithmetic hypothesis is a stratified Diophantine
bound
\[
 |a_j(k)|\ge c_j(1+|k|)^{-\tau_j}\qquad(k\in\Sigma_j).
\]
For a term of length $n$ in~\eqref{eq:mult-form}, its multiplier then grows
at most polynomially with exponent $n+(n+1)\tau_j$.
Thus every individual coefficient operator has explicit finite loss
estimates between narrower-strip seminorms.  There is no assertion that
these losses remain bounded as $\eta$ varies.  The fixed-strip theorem is
possible precisely because the algebra requires a common open domain, not
a single bounded analytic norm for every coefficient.

\begin{corollary}[Higher tails do not change the arithmetic test]\label{cor:tail}
If $\widetilde\Omega-\Omega$ has valuation greater than $\kappa$, then it
has the same flag, the same leading real forms, and the same exact resonance
lattice whenever $\res=\{0\}$.  In that nonresonant case,
$\widetilde\Omega$ satisfies $(\SDcond)$ exactly when $\Omega$ does.
\end{corollary}
\begin{proof}
No coefficient at or below $\kappa$ changes.  Each nonzero integer vector
already has a nonvanishing leading divisor coefficient at one of those
levels, so no resonance can appear and its leading form is unchanged.
\end{proof}

The nonresonance qualification matters.  When $\res$ is nontrivial, a new
higher tail can split some of its exact resonances and introduce additional
flag levels.  The criterion must then be applied to the extended flag.

\section{Arbitrarily long analytic jets with no analytic lift}
\label{sec:jets}

We now give a concrete failure that a valuation-only argument cannot detect.
All exponents in this section lie in $\Z$, and $t$ is a positive
infinitesimal.  The obstruction is therefore not caused by exotic
uncountable supports.

Define integers and a real number by
\begin{equation}\label{eq:alpha}
 b_1=2,\qquad b_{j+1}=10^{2b_j},\qquad
 \alpha=\sum_{j\ge1}10^{-b_j}.
\end{equation}
Put
\begin{equation}\label{eq:approximants}
 q_j=10^{b_j},\qquad
 p_j=q_j\sum_{\ell\le j}10^{-b_\ell},\qquad
 a_j=q_j\alpha-p_j,\qquad
 k_j=(-p_j,q_j),\quad K_j=|k_j|.
\end{equation}
Here $p_j$ is an integer, $0<p_j<q_j$, and
\begin{equation}\label{eq:liouville}
 0<a_j\le2q_j10^{-q_j^2},\qquad q_j\le K_j<2q_j.
\end{equation}
The error estimate follows by bounding the tail of the rapidly increasing
sequence of decimal places.  It also proves irrationality: if $\alpha$ had
fixed denominator $q$, each nonzero $a_j$ would have modulus at least $1/q$,
contrary to~\eqref{eq:liouville}.

Consider the infinitesimally detuned frequency and operator
\begin{equation}\label{eq:detuned}
 \Omega=(1,\alpha+t),\qquad
 \Lom=\partial_x+(\alpha+t)\partial_y=L_0+t\partial_y,
 \quad L_0=\partial_x+\alpha\partial_y.
\end{equation}
The real leading vector $(1,\alpha)$ is already nonresonant, so there is
one flag level, at zero.  The small-divisor condition fails very strongly.

\begin{theorem}[Arbitrary finite analytic lifting, but no full lift]
\label{thm:jet}
For every integer $N\ge1$ and every $\rho>0$, there exists an ordinary real
entire periodic function $f_N$ such that:
\begin{enumerate}[label=(\roman*),leftmargin=*]
\item The equation $\Lom u=f_N$ has a mean-zero solution modulo $t^{N+1}$
with coefficients $u_0,\ldots,u_N\in\A_\rho$.
\item The coefficients $u_0,\ldots,u_{N-1}$ are entire periodic functions.
\item The forced coefficient $u_{N+1}$ is not a distribution on the real
torus.  Consequently there is no solution in $\HH_{\rho'}$ for any
$\rho'>0$, even after an order-preserving extension of the value group.
\end{enumerate}
One may take
\begin{equation}\label{eq:jet-forcing}
 f_N(x,y)=2\sum_{j\ge1}a_j^{N+1}e^{-\rho K_j}
                                  \cos(k_j\cdot(x,y)).
\end{equation}
\end{theorem}
\begin{proof}
Estimate~\eqref{eq:liouville} shows that the Fourier coefficients
in~\eqref{eq:jet-forcing} decay faster than $e^{-R K_j}$ for every fixed
$R>0$.  Hence $f_N$ is entire.  It is real and has zero mean.

At the mode $k_j$, the divisor is $a_j+q_jt$.  Expanding its inverse as a
Hahn geometric series gives the forced coefficients
\begin{equation}\label{eq:jet-coefficients}
 \widehat{u_n}(k_j)=
 \frac{(-1)^n}{i}\,q_j^n a_j^{N-n}e^{-\rho K_j}
 \qquad(n\ge0).
\end{equation}
At $-k_j$ the coefficients are their complex conjugates; all other nonzero
modes vanish.  Set every mean to zero.

For $n<N$, the remaining positive power of $a_j$ gives superexponential
Fourier decay, so $u_n$ is entire.  For $n=N$, the coefficients are a
polynomial factor $q_j^N$ times $e^{-\rho K_j}$; their weighted Fourier sums
are finite on every strip of width $r<\rho$.  Thus $u_N\in\A_\rho$.
The scalar identities imply
\[
 L_0u_0=f_N,\qquad L_0u_n=-\partial_yu_{n-1}\quad(1\le n\le N),
\]
which prove the finite-jet statement in $\A_\rho[[t]]/(t^{N+1})$.

At the next coefficient,
\[
 |\widehat{u_{N+1}}(k_j)|
 =q_j^{N+1}a_j^{-1}e^{-\rho K_j}
 \ge\tfrac12 q_j^N10^{q_j^2}e^{-2\rho q_j}.
\]
This exceeds every polynomial in $K_j$ along the sequence.  It cannot be
the Fourier coefficient sequence of a distribution.

Finally, this obstruction cannot be evaded by allowing negative or
noninteger Hahn exponents in a proposed solution.  At each nonzero Fourier
mode the scalar equation has the unique solution
$\widehat f_N(k)/(i\,k\cdot\Omega)$ in the Hahn field.  Its inverse is still
the same geometric series in any larger ordered-group Hahn field.  Therefore
the coefficient of $t^{N+1}$ remains exactly~\eqref{eq:jet-coefficients}.
Extra constants affect only the zero Fourier mode and cannot repair it.
\end{proof}

This is a finite-versus-infinite lifting obstruction in a precise analytic
ring.  It does not say that the scalar algebraic inverses fail to exist;
they all exist.  It says that infinitely many modewise inverses do not
assemble into even one analytic coefficient at a specified order.

The same operator~\eqref{eq:detuned} works for all $N$; only the entire
forcing is changed.  In particular, no test consisting of a fixed finite
number of analytic coefficient equations can replace the arithmetic
hypothesis in a general lifting theorem.  The construction adapts the
classical Liouville mechanism~\cite{kamen}, but makes its failure occur at
an arbitrarily prescribed perturbation order.

\section{Infinitesimal conjugacy with a frequency correction}
\label{sec:nonlinear}

We turn from scalar equations to vector fields.  By multiplying the entire
vector field by a scalar monomial, we may normalize
\begin{equation}\label{eq:normalized}
 \gamma_1=0,\qquad 0\le\kappa=\gamma_r.
\end{equation}
This normalization subtracts $\gamma_1$ from both frequency and perturbation
valuations, so the eventual condition $v(F)>\kappa$ is unchanged when translated
back to the original scale.
Assume $\Gamma$ is nontrivial, $\res=\{0\}$, and $(\SDcond)$ holds.
Write $\mathsf Q f=f-\mean f$ for the mean-zero projection.  The solution
operator $\Gom$ now inverts $\Lom$ on all mean-zero functions, and is applied
componentwise to vector-valued functions.

Let
\begin{equation}\label{eq:vectorfield}
 X(\theta)=\Omega+F(\theta),\qquad F\in\HH_\rho^d,
 \quad s_0=v(F)>\kappa.
\end{equation}
A near-identity coordinate transformation means
$\Phi=\id+h$, where $h$ is periodic and has positive valuation.  Composition
is Taylor composition from Lemma~\ref{lem:taylor}.

\begin{theorem}[Hahn-analytic constant normal form]\label{thm:nonlinear}
Under the preceding assumptions there is a unique pair
\[
 h\in\HH_\rho^d,\qquad a\in K_\Gamma^d,
 \qquad v(h)>0,\qquad\mean h=0,
\]
such that
\begin{equation}\label{eq:conjugacy}
 D\Phi(\theta)(\Omega+a)=X(\Phi(\theta)),\qquad \Phi=\id+h.
\end{equation}
The following additional conclusions hold.
\begin{enumerate}[label=(\roman*),leftmargin=*]
\item $v(h)\ge s_0-\kappa$ and $v(a)\ge s_0$.
\item If $B=\supp F$, set
\begin{equation}\label{eq:nonlinear-monoid}
 P_0=\bigcup_{j=1}^r(B-\gamma_j),\qquad
 P=E\cup P_0,\qquad M=\mon P.
\end{equation}
These are positive generators in a well-ordered set, and
\[
 \supp h\subseteq M\setminus\{0\},\qquad
 \supp a\subseteq B+M.
\]
\item $\Phi$ has an infinitesimal Taylor inverse on $\HH_\rho$.
All nonzero integer divisors of $\Omega+a$ have the same leading real
coefficients and valuations as those of $\Omega$.
\item If $F$ has real values coefficientwise on the real torus, then $h$
and $a$ are real in the same sense, with $a\in F_\Gamma^d$.
\end{enumerate}
The assertion also includes $F=0$, with $h=a=0$ and the evident empty-support
conventions.
\end{theorem}

\subsection{Reduction to a fixed-point equation}

Taking the mean of~\eqref{eq:conjugacy} uses $\mean{Dh}=0$ and yields
\begin{equation}\label{eq:drift}
 a=a(h):=\mean{F\circ(\id+h)}.
\end{equation}
The remaining equation becomes
\begin{equation}\label{eq:fixed-point}
 h=\mathcal T(h):=
 \Gom\left(\mathsf Q(F\circ(\id+h))-Dh\,a(h)\right).
\end{equation}
The argument of $\Gom$ has zero mean: $a(h)$ is constant and every column
of $Dh$ has zero mean.  Thus the inverse is used only where it is defined.

All terms in $B-\gamma_j$ are positive because $s_0>\kappa\ge\gamma_j$.
The sets in~\eqref{eq:nonlinear-monoid} are well ordered, and
Lemma~\ref{lem:neumann} applies.  If $h$ has positive support in $M$, then
Taylor substitution gives
\[
 \supp(F\circ(\id+h))\subseteq B+M,
 \qquad \supp a(h)\subseteq B+M.
\]
The same is true of $Dh\,a(h)$ after adding $M$ once more.  Applying an
inverse coefficient at $-\gamma_j+M_E$ shifts these supports into
$B-\gamma_j+M\subseteq M$.  Every resulting term includes at least one
letter from $P_0$, so no constant exponent is introduced.

This shows preservation of a support certificate.  Existence still requires
proof that infinitely many iteration steps assemble into a Hahn series.

\subsection{Existence by marked support words}

\begin{proof}[Existence in Theorem~\ref{thm:nonlinear}]
Start with $h^{(0)}=0$ and define $h^{(n+1)}=\mathcal T(h^{(n)})$.
Each iterate lies in $\HH_\rho^d$, has zero mean, and has positive support in
$M$.  Its coefficients are obtained by a finite number of analytic
operations at each exponent.

For the stabilization argument, regard a generator in $P_0$ as a
\emph{marked} letter and a generator in $E$ as unmarked.  If an exponent
belongs to both sets, keep both labels.  There are only finitely many
origins for a given letter: $j$ has finitely many possibilities, and
translation determines the original exponent of $B$ or $S$.  Hence the
finite-word conclusion of Lemma~\ref{lem:neumann} is unchanged by these
labels.

A term of $\mathcal T(h)$ consists of a coefficient of $F$, an inverse
coefficient, and finitely many factors of $h$ or its derivatives.  The
combined exponent of the $F$ coefficient and the leading shift of that
inverse supplies one marked letter $\beta-\gamma_j$.  The higher inverse
terms supply unmarked letters in $E$.  All other letters come from the
factors of $h$.

More precisely, by induction the difference
\begin{equation}\label{eq:iterate-differences}
 h^{(n+1)}-h^{(n)}
\end{equation}
can have a nonzero coefficient only at an exponent represented by a word
with at least $n+1$ marked letters.  For $n=0$ there is at least one $F$
factor.  For the induction step, subtract the two Taylor expressions and
telescope the finite products at each coefficient.  Every difference term
contains a factor from $h^{(n)}-h^{(n-1)}$ or its derivative, together with
one additional $F$ coefficient.  The term $Dh\,a(h)$ has the same property:
subtract it as $D(h-g)a(h)+Dg(a(h)-a(g))$, and expand the second difference
by the same Taylor telescoping.  Applying $\Gom$ creates one new marked
shift and unmarked inverse-tail letters, exactly as above.

For a fixed exponent $\gamma$, Lemma~\ref{lem:neumann} permits only finitely
many labelled word representations of $\gamma$.  In particular their
numbers of marked letters are bounded.  Thus~\eqref{eq:iterate-differences}
has zero coefficient at $\gamma$ for all sufficiently large $n$.
The coefficients $[t^\gamma]h^{(n)}$ stabilize exactly, not merely
approximately.  The union of the difference supports is contained in the
well-ordered set $M$, and each exponent occurs in only finitely many
differences.  Therefore these differences form a Hahn-summable family.
Let $h$ be their Hahn sum.

Each stabilized coefficient is a member of $\A_\rho$; no new analytic
limit has been taken.  The expression $\mathcal T(h)$ at a fixed exponent
uses only finitely many of the certified word decompositions and finitely
many stabilized input coefficients.  Passing to those stabilized
coefficients in the iteration proves $h=\mathcal T(h)$.
Define $a$ by~\eqref{eq:drift}.  Equations~\eqref{eq:fixed-point}
and~\eqref{eq:drift} give~\eqref{eq:conjugacy}.

Since substitution does not lower valuation, $v(a)\ge s_0$.
The right side inside $\Gom$ has valuation at least $s_0$, because
$Dh$ has positive valuation.  The linear valuation estimate gives
$v(h)\ge s_0-\kappa$.  The construction already proves the support
statements.
\end{proof}

The word argument is the essential existence step.  The weaker estimate
$v(h^{(n+1)}-h^{(n)})\ge(n+1)(s_0-\kappa)$ would not suffice in a general
ordered group: the sequence of those lower bounds need not be cofinal.
The proof does not use a Banach fixed-point theorem or sequential
completeness of the full surreal field.

\subsection{Uniqueness and the inverse coordinate map}

\begin{proof}[Uniqueness in Theorem~\ref{thm:nonlinear}]
Let $h,g$ be two positive-valuation, mean-zero solutions.  Their drifts are
forced by~\eqref{eq:drift}.  If $h\ne g$, put $\delta=v(h-g)>0$.
Lemma~\ref{lem:taylor} gives
\[
 v(a(h)-a(g))\ge s_0+\delta.
\]
Subtracting the expressions inside~\eqref{eq:fixed-point} gives valuation
at least $s_0+\delta$, since $a(h)$ has valuation at least $s_0$ and
differentiation does not lower valuation.  Hence
\[
 v(\mathcal T(h)-\mathcal T(g))
 \ge s_0+\delta-\kappa>\delta.
\]
This contradicts $h-g=\mathcal T(h)-\mathcal T(g)$.  Uniqueness holds among
all positive-valuation solutions, not just those carrying the particular
support certificate constructed above.
\end{proof}

\begin{lemma}[Infinitesimal coordinate inverses]\label{lem:coordinate-inverse}
For every $h\in\HH_\rho^d$ of positive valuation, $\Phi=\id+h$ has a
unique inverse $\Psi=\id+g$ of positive valuation in the Taylor sense.
It induces an automorphism of $\HH_\rho$.
\end{lemma}
\begin{proof}
Solve $g=-h\circ(\id+g)$ by iteration from zero.  A word argument identical
to the preceding existence proof, now marking each occurrence of a
coefficient of $h$, gives coefficient stabilization inside the positive
monoid generated by $\supp h$.  This constructs $\Phi\circ\Psi=\id$.
Uniqueness follows from
\[
 v\bigl(h\circ(\id+g_1)-h\circ(\id+g_2)\bigr)
 \ge v(h)+v(g_1-g_2)>v(g_1-g_2).
\]
To see that the inverse is two-sided, write $\Psi\circ\Phi=\id+q$.
Associativity gives $\Phi\circ(\id+q)=\Phi$, hence
$q+h\circ(\id+q)-h=0$.  If $q\ne0$, the second difference has valuation
strictly greater than $v(q)$, an impossibility.  The Taylor identities now
show that the two substitution operators are inverse ring homomorphisms.
\end{proof}

\begin{proof}[Completion of Theorem~\ref{thm:nonlinear}]
Lemma~\ref{lem:coordinate-inverse} proves the inverse assertion.
Since $v(a)\ge s_0>\kappa$, no leading coefficient of a nonzero integer
divisor changes.  This gives the flag conclusion of Corollary~\ref{cor:tail}
for real data; for complex data the same leading real coefficients persist,
although some higher coefficients can be complex.  Complex conjugation together with $k\mapsto-k$
preserves all scalar Fourier inversions for a real frequency.  Taylor
substitution, mean, and differentiation preserve coefficientwise reality.
Thus the iteration constructs real $h$ and $a$ from real $F$; alternatively,
reality follows directly from uniqueness.
\end{proof}

\section{Necessity for non-linearization and the sharp scale boundary}
\label{sec:necessity}

The non-linear hypothesis is not merely a convenient sufficient condition.
Within the specified Hahn-analytic category, the same arithmetic condition
is exact for a universal theorem.

\begin{theorem}[Arbitrarily small non-linearization obstructions]
\label{thm:nonlinear-necessity}
Assume $\gamma_1=0$, $\res=\{0\}$, and $\Gamma$ is nontrivial.  If
$(\SDcond)$ fails, then for every prescribed $H\in\Gamma$ there are $s>H$,
a real $\rho>0$, and a real perturbation $F\in\HH_\rho^d$ with $v(F)=s$
such that $X=\Omega+F$ admits no positive-valuation, mean-zero
Hahn-analytic conjugacy to a constant vector field, even on a narrower
positive-width strip.  One may require $s>2\kappa$ and take $F$ to have
just one Hahn exponent.
\end{theorem}
\begin{proof}
Choose $s>\max\{H,2\kappa\}$, which is possible in a nontrivial ordered
abelian group.  Use the real function $f\in\A_\eps$ from
\eqref{eq:bad-forcing} on a failing stratum $\Sigma_j$, and put
$F=t^s f e_1$, where $e_1$ is the first coordinate vector.

Suppose $\Phi=\id+h$ were such a conjugacy on some strip.  Restrict to a
common narrower strip if necessary.  Taking means still forces
$a=\mean{F\circ\Phi}$ and $v(a)\ge s$.  Since $v(h)>0$, the equation
\[
 \Lom h=\mathsf Q(F\circ\Phi)-Dh\,a
\]
has a right side of valuation at least $s$.  The arithmetic-free estimate
in Proposition~\ref{prop:val-bound} gives $v(h)\ge s-\kappa$.

Subtract the uncomposed forcing to write
\[
 \Lom h=t^s f e_1+R,
 \qquad
 R=\mathsf Q(F\circ\Phi-F)-Dh\,a.
\]
Taylor telescoping gives
\[
 v(R)\ge s+v(h)\ge2s-\kappa.
\]
For every nonzero Fourier mode the scalar equation therefore implies
\[
 \widehat h(k)=\frac{t^s\widehat f(k)e_1}{i\,k\cdot\Omega}
               +\frac{\widehat R(k)}{i\,k\cdot\Omega},
 \qquad
 v\left(\frac{\widehat R(k)}{i\,k\cdot\Omega}\right)
       \ge2s-2\kappa>s.
\]
For $k=k_n$ in the chosen failing stratum, the coefficient of
$t^{s-\gamma_j}$ in the first component of $h$ is consequently forced to
be $e^{-\eps|k_n|}/(i a_j(k_n))$.  Here $s-\gamma_j\le s$, so the
remainder cannot contribute.  These coefficients grow faster than
$e^{\eps|k_n|}$ and cannot belong to a distribution, exactly as in the
necessity proof of Theorem~\ref{thm:linear}.  This contradiction proves the
assertion.
\end{proof}

\begin{corollary}[Exact universal infinitesimal normal-form criterion]
\label{cor:nonlinear-equivalence}
For a normalized nonresonant $\Omega$ over a nontrivial $\Gamma$, condition
$(\SDcond)$ is equivalent to the assertion that, for every $\rho>0$, every
$F\in\HH_\rho^d$ of valuation greater than $\kappa$ admits the normalized
conjugacy of Theorem~\ref{thm:nonlinear}.  When $(\SDcond)$ fails, restricting
perturbations to arbitrarily higher valuations does not rescue a universal
theorem.
\end{corollary}

\subsection{Why strict inequality at the final scale is necessary}

\begin{proposition}[Failure at the threshold, without a zero of the slow speed]
\label{prop:boundary}
Let $\kappa>0$ belong to $\Gamma$, and consider
\[
 \Omega=(1,t^\kappa),\qquad
 F(\theta_1,\theta_2)=\left(0,\tfrac12t^\kappa\cos\theta_2\right).
\]
The frequency satisfies $(\SDcond)$ and its final flag level is $\kappa$.
The perturbation has valuation $\kappa$, but $\Omega+F$ has no
positive-valuation, mean-zero conjugacy to a constant vector field.
\end{proposition}
\begin{proof}
The two strata have leading forms $k_1$ and $k_2$, respectively, so
$(\SDcond)$ is immediate.  Suppose a conjugacy existed.  The mean equation
forces $a_1=0$.  Since substitution by a positive-valuation $h$ does not
change the coefficient at exponent $\kappa$ of $F$, that coefficient is
$(0,\tfrac12\cos\theta_2)$, of mean zero.  Hence $v(a_2)>\kappa$, with
$a_2=0$ allowed.

At exponent $\kappa$ in the second component of the conjugacy equation,
$t^\kappa\partial_2h_2$ has no contribution because $[t^0]h_2=0$.
The drift and $Dh\,a$ have no contribution either.  What remains is
\[
 \partial_1([t^\kappa]h_2)=\tfrac12\cos\theta_2.
\]
Averaging in $\theta_1$ makes the left side zero and leaves the nonzero
right side unchanged, a contradiction.
\end{proof}

The slow component of this vector field is
$t^\kappa(1+\tfrac12\cos\theta_2)$, whose leading real factor is everywhere
positive.  Thus the example does not rely on introducing a stationary point.
It says that a perturbation of the \emph{same} size as a slow frequency may
require a coordinate correction of order one, rather than an infinitesimal
correction.  The strict threshold in the theorem is therefore sharp as a
uniform sufficient bound; particular perturbations at the boundary can of
course be linearizable.

\section{A canonical invariant density}
\label{sec:density}

The conjugacy gives more than a constant normal form.  It yields a canonical
Hahn-valued averaging functional and an analytic density for it.  The claim
concerns coefficientwise integration of analytic functions, not a
countably additive $\No$-valued probability measure on arbitrary subsets.

\begin{lemma}[Formal change of variables]\label{lem:change-variables}
For $\Phi=\id+h$ with $v(h)>0$, and every $f\in\HH_\rho$,
\begin{equation}\label{eq:change-variables}
 \mean{(f\circ\Phi)\det D\Phi}=\mean f.
\end{equation}
\end{lemma}
\begin{proof}
Introduce an ordinary formal auxiliary variable $z$ and
$\Phi_z=\id+zh$.  At each Hahn exponent, Taylor substitution and the
determinant produce a polynomial in $z$: there are only finitely many
support words at that exponent.  Put $A=D\Phi_z$.  The finite determinant
identities and the chain rule give
\begin{equation}\label{eq:piola}
 \frac{\mathrm d}{\mathrm dz}
       \bigl((f\circ\Phi_z)\det A\bigr)
 =\diver_\theta\bigl((f\circ\Phi_z)\operatorname{adj}(A)h\bigr).
\end{equation}
To check the identity, use $A\operatorname{adj}(A)=(\det A)I$ for the
derivative of $f$, and the derivative-of-determinant formula for the terms
involving $Dh$.  The remaining derivatives of the cofactor matrix cancel
in pairs because mixed partial derivatives of $\Phi_z$ commute; this is
the elementary Piola identity for a Jacobian matrix.

Each coefficient on the right of~\eqref{eq:piola} is a divergence of a
periodic analytic vector field and has mean zero.  Therefore the mean on
the left is independent of $z$.  Evaluating the coefficient polynomials at
$z=0$ and $z=1$ gives~\eqref{eq:change-variables}.  There is no infinite
integration in $z$ and no assumption of ordinary convergence of a Hahn sum.
\end{proof}

\begin{theorem}[Invariant density and its uniqueness]\label{thm:density}
Let $X$, $\Phi$, and $a$ be as in Theorem~\ref{thm:nonlinear}, and put
$\Psi=\Phi^{-1}$.  Then
\begin{equation}\label{eq:density}
 \mu_X=\det D\Psi\in 1+\HH_{\rho,>0}
\end{equation}
is the unique density in $\HH_\rho$ satisfying
\begin{equation}\label{eq:density-properties}
 \diver(\mu_X X)=0,\qquad \mean{\mu_X}=1.
\end{equation}
The $K_\Gamma$-linear functional
\begin{equation}\label{eq:invariant-functional}
 I_X(f)=\mean{f\circ\Phi}=\mean{f\mu_X}
\end{equation}
satisfies $I_X(1)=1$ and $I_X(X\cdot\nabla f)=0$.
For real data, $\mu_X$ is real and has positive surreal values on the real
torus, since it is one plus an infinitesimal.
\end{theorem}
\begin{proof}
The inverse is infinitesimal, so its Jacobian determinant is one plus an
element of positive valuation.  The chain rule gives
\[
 (\mu_X\circ\Phi)\det D\Phi=1.
\]
Apply Lemma~\ref{lem:change-variables} to $f\mu_X$ to obtain the equality
in~\eqref{eq:invariant-functional}.  With $f=1$ it gives normalization.
Conjugacy and the chain rule yield
\[
 (X\cdot\nabla f)\circ\Phi
       =(\Omega+a)\cdot\nabla(f\circ\Phi),
\]
whose mean is zero.  Thus $I_X(X\cdot\nabla f)=0$.
Coefficientwise integration by parts, tested against every Fourier
character, implies $\diver(\mu_X X)=0$.

For uniqueness, let $\widetilde\mu$ be any normalized solution of
\eqref{eq:density-properties} and set
\[
 g=(\widetilde\mu\circ\Phi)\det D\Phi.
\]
The same Jacobian identities transport its divergence equation to
$(\Omega+a)\cdot\nabla g=0$.  Since $\Omega+a$ has no nonzero integer
resonances, Fourier uniqueness implies that $g$ is constant in $\theta$.
The change-of-variables identity gives
$\mean g=\mean{\widetilde\mu}=1$, so $g=1$ and
$\widetilde\mu=\mu_X$.  Reality and pointwise positivity follow from the
reality assertion in Theorem~\ref{thm:nonlinear} and the positive-valuation
form of~\eqref{eq:density}.
\end{proof}

This result is an algebraic-analytic invariant-density statement.  It does
not assert existence of a real-time flow taking values continuously in the
fine topology of all surreals, nor an ergodic theorem involving such a flow.
The identity $I_X(X\cdot\nabla f)=0$ is the precise invariant assertion.

\section{Examples with several genuinely different scales}
\label{sec:examples}

\subsection{Three visible scales and a support of order type \texorpdfstring{$\omega^2$}{omega squared}}

Take $\Gamma=\Z+\Z\omega\subset\No$ and
\begin{equation}\label{eq:three-frequencies}
 \Omega=(1,t,t^\omega).
\end{equation}
Here the symbol $\omega$ in the exponent is the positive infinite surreal
number.  The flag is
\[
 \Z^3\supset\{k_1=0\}\supset\{k_1=k_2=0\}\supset\{0\},
\]
with levels $0,1,\omega$ and leading forms $k_1,k_2,k_3$ on their respective
strata.  Each nonzero leading form is a nonzero integer, so its reciprocal
has modulus at most one.  Thus $(\SDcond)$ holds without any delicate ordinary
arithmetic.

Already the mode $k=(1,1,1)$ has the inverse
\begin{equation}\label{eq:ranktwo-inverse}
 \frac1{1+t+t^\omega}
 =\sum_{n,m\ge0}(-1)^{n+m}\binom{n+m}{n}\,t^{n\omega+m}.
\end{equation}
The support has order type $\omega^2$: every exponent in the $n=0$ layer
precedes every exponent in the $n=1$ layer, and so on.  Formula
\eqref{eq:ranktwo-inverse} follows by grouping words according to the numbers
of letters $1$ and $\omega$.  It is also verified by the finite coefficient
recurrence
\[
 C_{n,m}+C_{n-1,m}+C_{n,m-1}
 =\begin{cases}1,&n=m=0,\\0,&\text{otherwise},\end{cases}
\]
where $C_{n,m}=(-1)^{n+m}\binom{n+m}{n}$ and a negative index means zero.

For a mode on the second stratum,
\[
 \frac1{k_2t+k_3t^\omega}
 =\frac{t^{-1}}{k_2}
  \sum_{n\ge0}\left(-\frac{k_3}{k_2}\right)^n t^{n(\omega-1)}.
\]
For a mode on the last stratum, the inverse is simply $t^{-\omega}/k_3$.
Thus the large support in~\eqref{eq:ranktwo-inverse} coexists with a finite
three-level divisor classification.

The scalar equation with
\[
 f=t^{\omega+1}(\sin\theta_1+\sin\theta_2+\sin\theta_3)
\]
has the exact mean-zero solution
\begin{equation}\label{eq:three-solution}
 u=-t^{\omega+1}\cos\theta_1-t^\omega\cos\theta_2-t\cos\theta_3.
\end{equation}
Its valuation is $1=v(f)-\omega$, so the last scale really controls the
worst inverse loss.  A vector perturbation of valuation $\omega+1$ likewise
falls under the non-linear theorem, with a coordinate correction of
valuation at least one.

This example exposes the failure of a sequential contraction proof in a
particularly simple setting.  Repeated corrections of order $1,2,3,\ldots$
need not pass the exponent $\omega$.  The coefficient-stabilization proof
still constructs the exact transformation, including coefficients at and
beyond that scale.

\subsection{Real arithmetic on a leading stratum}

Take $\Omega=(1,\sqrt2,t)$.  Its first stratum consists of all
$(p,q,r)$ with $(p,q)\ne(0,0)$ and has leading form $p+\sqrt2q$.
Its second stratum is $(0,0,r)$ with $r\ne0$, at level one.  Since
\[
 |p+\sqrt2q|^{-1}
 =\frac{|p-\sqrt2q|}{|p^2-2q^2|}
 \le |p|+\sqrt2|q|
\]
for $(p,q)\ne(0,0)$, the first stratum is Diophantine.  The second has the
integer leading form $r$.  The universal conjugacy threshold is therefore
$v(F)>1$.

One may append an arbitrary set-supported real vector tail of valuation
greater than one.  It can have transfinite support and change every higher
coefficient of the solution, but it does not change this arithmetic test or
the threshold.  This is a useful separation of ordinary approximation
properties from surreal support complexity.

\subsection{An explicit slow-circle normal form}

Let $q=t^\kappa$, $b=t^s$, where $s>\kappa>0$, and consider
\[
 X(\theta_1,\theta_2)=(1,q+b\sin\theta_2).
\]
Set $r=b/q=t^{s-\kappa}$.  The mean-zero conjugacy has the form
$\Phi=(\theta_1,\theta_2+h(\theta_2))$ and transforms $X$ to $(1,\nu)$,
where
\begin{equation}\label{eq:circle-conjugacy}
 (1+h'(\theta))\nu=q+b\sin(\theta+h(\theta)).
\end{equation}
The new frequency is exactly
\begin{equation}\label{eq:circle-frequency}
 \nu=q\sqrt{1-r^2}
 =q-\frac{b^2}{2q}-\frac{b^4}{8q^3}-\frac{b^6}{16q^5}-\cdots,
\end{equation}
where the square root is the Hahn binomial branch with constant term one.

To prove the formula, use the inverse coordinate $\Psi$ in the second
variable.  The inverse equation gives
\[
 \Psi'(\theta)=\frac{\nu}{q(1+r\sin\theta)}.
\]
Its mean is one, so
\[
 \frac q\nu=\mean{(1+r\sin\theta)^{-1}}
   =\sum_{m\ge0}\binom{2m}{m}\frac{r^{2m}}{4^m}
   =(1-r^2)^{-1/2}.
\]
All identities are Hahn identities in the positive monomial $r$.
The middle equality follows from the geometric expansion, the vanishing
of odd sine moments, and
$\mean{\sin^{2m}\theta}=\binom{2m}{m}/4^m$, obtained by ordinary integration
by parts or Fourier coefficient extraction.  The binomial identity proves
\eqref{eq:circle-frequency}.

The first terms of the unique mean-zero $h$ are
\begin{equation}\label{eq:circle-h}
 h(\theta)=-r\cos\theta-\frac{r^2}{4}\sin2\theta
       +r^3\left(\frac1{12}\cos3\theta-\frac14\cos\theta\right)
       +O_H(r^4).
\end{equation}
Here $O_H(r^4)$ means the remaining integer powers of $r$ have exponent at
least four; it is not an assertion about convergence in a real parameter.
In particular, the lower bound $v(h)\ge s-\kappa$ is attained.  The first
frequency drift occurs later, at valuation $2s-\kappa$.

The invariant density is equally explicit:
\begin{equation}\label{eq:circle-density}
 \mu_X(\theta_2)=\frac{\sqrt{1-r^2}}{1+r\sin\theta_2}.
\end{equation}
It has mean one and makes $\mu_X(q+b\sin\theta_2)$ independent of
$\theta_2$.  The verification program checks~\eqref{eq:circle-conjugacy}
through degree ten in $r$, independently reconstructing $h$ and $\nu$ by
exact Fourier arithmetic.

\section{Interpretation, novelty boundaries, and unresolved extensions}
\label{sec:limits}

\subsection{What the theorem says about surreal analysis}

The full surreal field permits many incomparable asymptotic rates, but a
finite vector cannot expose an unlimited chain of integer resonance
kernels: their ranks decrease.  This gives a finite arithmetic skeleton
inside a transfinite coefficient expansion.  Positive-support combinatorics
then turns that skeleton into analytic inverse operators and exact
non-linear normal forms.

The resulting objects can be evaluated at infinitesimal thickenings of
ordinary torus points by~\eqref{eq:monad-eval}.  In that precise sense these
are statements about surreal- and surcomplex-valued functions, not merely
statements about an unrelated formal indeterminate.  Yet the proofs
localize to set-sized fields, so they do not need a class-sized topology or
class-indexed analytic limits.

The rank bound alone is an elementary finite-dimensional observation.
Its significance here is that it yields \emph{uniform control of infinitely
many scalar inverses}, including their common Hahn support.  It is not
presented as a stand-alone major discovery.

\subsection{Why this is not an improvement of real KAM arithmetic}

A formal series $\sum_{n\ge0}t^n h_n(\theta)$ is an element of $\HH_\rho$
whenever every $h_n$ is holomorphic on the same strip, regardless of how
rapidly $\|h_n\|_r$ grows.  Even growth faster than every exponential in $n$
is permitted.  Therefore substitution of a nonzero real number for $t$
need not be meaningful.

Classical analytic linearization is a different problem: it must control
both Fourier decay and accumulation of losses through infinitely many
non-linear steps.  Brjuno-type conditions enter such convergence theorems
\cite{bounemoura,chevallier}.  Our use of $(\SDcond)$ only guarantees that each
\emph{individual} coefficient operation remains holomorphic; the Hahn
support argument handles how those coefficients are assembled.  No bound
uniform in the Hahn exponent is supplied.

Likewise, increasing valuations do not automatically give convergence even
inside a fixed higher-rank valued field.  The marked-word proof gives Hahn
summability, which is the exact operation required by the theorem.  These
two qualifications prevent a false transfer to either ordinary analytic KAM
theory or the fine topology on all surcomplex numbers.

\subsection{A conservative novelty audit}

The literature review included the named sources in the bibliography and
focused searches for Hahn series, surreal frequencies, small divisors,
cohomological equations, and torus linearization.  The exact theorem package
was not located.  Such searches do not establish absence from the entire
published literature, including work under different terminology.

\begin{center}
\small
\begin{tabular}{@{}p{0.47\textwidth}p{0.45\textwidth}@{}}
\toprule
Result or ingredient & Attribution/status in this article\\
\midrule
Positive-support inverses and finite word representations & Classical Hahn--Neumann input; not a novelty claim.\\[4pt]
Ordinary Fourier division and subexponential arithmetic & Classical small-divisor input, with an elementary proof included.\\[4pt]
Finite resonance flag plus one support for all reciprocal divisors & Proposed surreal/Hahn formulation and uniform-support theorem.\\[4pt]
Exact common-strip solvability across all strata & Proposed assembly theorem; depends on the preceding two classical tools.\\[4pt]
Arbitrarily delayed analytic lifting obstruction & Explicit construction proposed here; its Liouville mechanism is classical.\\[4pt]
Transfinite-support normal form, necessity, and sharp threshold & Principal proposed theorem package; no claim of real-parameter convergence.\\[4pt]
Invariant density and averaging & Derived consequence; the finite Jacobian identities are classical.\\
\bottomrule
\end{tabular}
\end{center}

No item is described as the resolution of an externally named open problem.
The contribution offered is a coherent group of proved statements in a
specified surcomplex analytic category.  Independent mathematical review is
still needed before treating their correctness and originality as settled.

\subsection{Three extensions that need additional work}

\paragraph{Quantitative coefficient-growth classes.}
One can add bounds that link $\|f_\beta\|_r$ to a weight on the support.
The present argument does not prove preservation of such a weighted class,
because inverse denominator powers and derivative losses grow along support
words.  An interesting next problem is to identify hypotheses that preserve
these weights and permit a meaningful specialization of selected monomials
to small real parameters.  That is where genuine non-linear convergence
arithmetic must re-enter.

\paragraph{Nontrivial exact resonances.}
The linear theorem allows $\res\ne\{0\}$, but the constant non-linear
normal form theorem does not.  In the resonant case a natural target is a
normal form retaining Fourier modes in $\res$.  Products and substitutions
interact with those modes, and higher perturbations may split the remaining
resonances.  Simply replacing the mean projection by $\Pi_{\res}$ in the
proof is not justified: the retained normal form is no longer a constant
vector and creates additional terms in the conjugacy equation.

\paragraph{Effective arithmetic certificates.}
The flag is finite, but not automatically computable from arbitrary real
coefficients.  Detecting an exact integer relation among arbitrary named
real numbers may itself be unavailable.  For rational and suitable
algebraic coefficient data, explicit lattice and approximation certificates
can make the construction effective.  The present theorems distinguish
existence of a finite flag from decidability of its entries and arithmetic
bounds.

These are directions arising from this article, not assertions that each
is a previously catalogued open problem in the literature.

\section{Verification record and a route to formalization}
\label{sec:verification}

\subsection{What was actually checked}

The accompanying program \texttt{code/verify.py} uses only the Python
standard library.  All numbers in its computations are exact elements of
$\Q(i)$ represented by pairs of rational numbers.  There is no floating
point threshold in any assertion.

For the slow-circle example, it reconstructs the mean-zero Fourier
coefficients recursively from
\[
 (1+h')\nu=1+r\sin(\theta+h),\qquad \nu(0)=1.
\]
At order $n$ it first determines the mean, hence the new coefficient of
$\nu$, and then integrates the remaining zero-mean Fourier polynomial.
It verifies the full conjugacy residual through degree ten and separately
checks $\nu^2=1-r^2$ through the same degree.  The recorded coefficients are
\[
 1,\ 0,\ -\tfrac12,\ 0,\ -\tfrac18,\ 0,\ -\tfrac1{16},\ 0,
 -\tfrac5{128},\ 0,\ -\tfrac7{256}.
\]
It also checks reality, zero mean, and the three terms displayed in
\eqref{eq:circle-h}.

For~\eqref{eq:ranktwo-inverse}, the program checks 117 exact coefficient
recurrences in a rectangular $(n,m)$ window.  It checks the three-stratum
classification for 1,330 nonzero integer modes in $[-5,5]^3$.  The stratum
counts are 1,210, 110, and 10.  All recorded assertions passed in the supplied
run; the machine-readable output is \texttt{verification.json}.

These computations do not verify irrationality of~\eqref{eq:alpha},
subexponential estimates at infinitely many modes, the word lemma, or the
infinite-support conjugacy theorem.  Those claims rest on the mathematical
proofs, not on the finite tests.  No claim of a Lean-checked or
Wolfram-verified infinite theorem is made.

\subsection{A plausible Lean development order}

The repository already exposes a Hahn support layer and an axiom-audited
finite algebra/foundations track, as described in its inspected root
README~\cite{repo}.  The proofs here suggest a modular extension rather
than a monolithic formalization.

The first component would formalize the finite flag in terms of subgroups
of $\Z^d$ and coefficient homomorphisms.  Its core claims are saturation,
strict rank decrease, and identification of the leading coefficient of a
nonzero divisor.  This part needs no analytic functions.

The second component would formalize the common reciprocal support and
coefficient formula using summable families.  Crucially, it should encode
\emph{a common support bound} separately from \emph{summability of a family};
the reciprocal family indexed by Fourier modes generally has only the
former property.

The third component would provide the fixed-strip Fourier algebra and its
subexponential multiplier lemma.  Once that bridge is in place, the linear
solution can be assembled by finite coefficient convolution.  The
non-linear component would then use a labelled positive-word relation,
with the number of marked letters as an auxiliary finite complexity bound
for each exponent.  Replacing this with a sequence of valuation estimates
would formalize the wrong argument.

Finally, finite formal differentiation and determinant identities support
the coordinate-inverse and density corollaries.  The negative examples
require a separate explicit Fourier-growth development.  No Lean source is
included here, because that complete analytic infrastructure was not built
and checked for this article.

\section{Conclusion}

Surreal small divisors have a finite resonance skeleton even when their
coefficients contain transfinite scale information.  That skeleton supplies
both the correct arithmetic condition and a uniform Hahn support for all
modewise inverses.  It turns ordinary Fourier division into a well-defined
common-strip surcomplex solution operator.

The non-linear application establishes an exact infinitesimal normal form
beyond the last resonance scale, with a unique coordinate correction, a
frequency drift, and a canonical invariant density.  Its converse shows
that valuation-smallness cannot overcome bad real arithmetic.  Its boundary
example shows that the relative scale threshold is essential.  Finally,
the arbitrary-order lifting construction demonstrates why any finite
collection of successful analytic coefficient calculations is insufficient
without a global arithmetic argument.

Together these results give a concrete answer to a structural question in
surreal analysis: transfinite support and analytic Fourier regularity can be
controlled simultaneously, but only when both are made explicit.  Neither
one replaces the other.

\clearpage
\appendix
\section{Proof-dependency and assumption audit}
\label{app:audit}

\begin{center}
\small
\begin{tabular}{@{}p{0.29\textwidth}p{0.32\textwidth}p{0.29\textwidth}@{}}
\toprule
Statement & Essential hypothesis & What fails without it\\
\midrule
Finite flag & Finite dimension $d$; Hahn scalar frequencies & Infinite-dimensional rank need not decrease to zero in finitely many steps.\\[4pt]
Common inverse support & One well-ordered support for the frequency vector & Independent unbounded scalar choices do not supply a shared support.\\[4pt]
Coefficient multiplier theorem & $(\SDcond)$ on every nonresonant stratum & A forced coefficient can fail to be a distribution.\\[4pt]
Fixed-strip solution & All input coefficients holomorphic on one open strip & Independent germs need not yield a common analytic domain.\\[4pt]
Non-linear constant form & No nonzero exact resonances & Resonant Fourier terms need not be removable.\\[4pt]
Positive-support construction & $v(F)>\kappa$ & The boundary example requires an order-one coordinate change.\\[4pt]
Non-linear existence proof & Finitely many support words at each exponent & Increasing valuations alone need not produce a Hahn sum.\\[4pt]
Uniqueness & Positive coordinate correction and mean-zero normalization & Constant translations remove the normalization; larger corrections lie outside the argument.\\[4pt]
Density interpretation & Coefficientwise torus mean and Jacobian calculus & No measure or real-time ergodic theorem follows automatically.\\
\bottomrule
\end{tabular}
\end{center}

\paragraph{Set size.}
Every support, Fourier index set, and coefficient family used in a proof is
set-sized.  The full surreal field enters only through embeddings of these
workspaces and the statement that a finite frequency vector can be
localized to one of them.

\paragraph{Order of summation.}
Analytic Fourier sums are taken inside coefficient functions.  Hahn sums
are assembled only after a common support and finite coefficient
multiplicity have been established.  The paper uses no unrestricted
interchange of these two summation notions.

\paragraph{Arithmetic versus valuation.}
The quantities $|a_j(k)|$ are ordinary real absolute values, while
$v(k\cdot\Omega)=\gamma_j$ lies in $\Gamma$.  Inequalities involving one
kind of size are never used as estimates for the other without an explicit
leading-coefficient formula.

\paragraph{Proof status.}
The finite checks are reproducible.  The proofs are mathematical arguments,
not proof-assistant certificates.  The proposed novelty assessment is
limited by the scope of the literature review and should be re-evaluated
before submission for publication.

\clearpage
\begin{thebibliography}{99}

\bibitem{gonshor}
H. Gonshor,
\emph{An Introduction to the Theory of Surreal Numbers},
London Mathematical Society Lecture Note Series 110,
Cambridge University Press, 1986.
\url{https://doi.org/10.1017/CBO9780511629143}.

\bibitem{neumann}
B. H. Neumann,
On ordered division rings,
\emph{Transactions of the American Mathematical Society} \textbf{66} (1949),
202--252.
\url{https://doi.org/10.1090/S0002-9947-1949-0032593-5}.

\bibitem{poonen}
B. Poonen,
\emph{Units in Hahn--Mal'cev--Neumann rings},
preprint dated January 14, 2026.
\url{https://math.mit.edu/~poonen/papers/malcev.pdf}.

\bibitem{bounemoura}
A. Bounemoura,
Optimal linearization of vector fields on the torus in non-analytic Gevrey
classes,
\emph{Annales de l'Institut Henri Poincar\'e C, Analyse Non Lin\'eaire}
\textbf{39} (2022), no.~3, 501--528.
\url{https://doi.org/10.4171/AIHPC/12}.
Preprint: \url{https://arxiv.org/abs/2006.05673}.

\bibitem{kamen}
P. Kamie\'nski,
\emph{The one-frequency cohomological equation, Brjuno-like functions and
Khintchine--L\'evy numbers},
preprint, 2019.
\url{https://arxiv.org/abs/1903.00311}.

\bibitem{chevallier}
N. Chevallier, J. Lopes Dias, J. P. Gaiv\~ao, A. Marnat, and N. Moshchevitin,
\emph{Brjuno condition through best approximations and the linearization
problem}, arXiv:2607.25610v1, July 2026.
\url{https://arxiv.org/abs/2607.25610}.

\bibitem{repo}
V. Reshetnikov, maintainer,
\emph{Surreal: Lean formalization and research reports},
repository root README and documentation catalogue, inspected
September 21, 2026.  Commit:
\begin{center}\small\texttt{aa846271b4dcae2c055b216126a87210292ec19b}\end{center}
\href{https://github.com/VladimirReshetnikov/Surreal/tree/aa846271b4dcae2c055b216126a87210292ec19b}{Pinned repository snapshot}.
The reviewed research reports are drafts; their presence in the repository
is not a claim of refereed or complete machine-checked status.

\end{thebibliography}
\end{document}
